\documentclass[11pt,a4paper]{article}
\usepackage[T2A]{fontenc}
\usepackage[utf8]{inputenc}
\usepackage[russian]{babel}
\usepackage{booktabs}
\usepackage{geometry}
\usepackage{amsmath}
\geometry{margin=2.2cm}

\title{Протокол байесовской предиктивной модели\\риска сейсмособытия (24 часа)}
\author{}
\date{\today}

\begin{document}
\maketitle

\section{Цель}
Построить модель вероятности сейсмособытия в следующие 24 часа:
\[
P(y_{t+1}=1 \mid \mathcal{F}_t),
\]
где признаки формируются только из данных до конца суток $t$.

\section{Реализация}
Подготовлен следующий пайплайн:
\begin{itemize}
  \item формирование ежедневного датасета:
  \texttt{seismic\_pipeline\_standalone/scripts/build\_seismic\_risk\_dataset.py};
  \item обучение и валидация байесовских моделей:
  \texttt{seismic\_pipeline\_standalone/scripts/run\_bayesian\_risk\_cv.py}.
\end{itemize}

\section{Разметка и контроль утечек}
\begin{enumerate}
  \item Ежедневные строки \texttt{rat\_id $\times$ sample\_date}; целевая
  переменная \texttt{label\_next24h=1}, если событие у этой крысы наступает
  на следующий день.
  \item Quiet controls отбираются из дней с буфером по расстоянию до любого
  события (\texttt{quiet\_buffer\_days}, по умолчанию 8 суток) и матчингом по
  крысе/месяцу.
  \item Артефакты исключаются до расчёта признаков:
  \texttt{row 0 day 1}, \texttt{row 4 day 3}, \texttt{row 5 day 3},
  \texttt{row 22 day 4}.
\end{enumerate}

\section{Feature engineering (causal)}
\begin{itemize}
  \item \textbf{Уровень}: mean/range/std за 24ч, 48ч, 72ч; первая и вторая
  половины суток.
  \item \textbf{Динамика}: тренды последних 3 дней, дельты к предыдущему дню,
  отношение mean текущего дня к среднему предыдущих 3 дней.
  \item \textbf{Форма}: \texttt{shape\_l1\_lag1}, \texttt{shape\_l2\_lag1},
  попарная дисперсия L1/L2 на последних 3 днях.
  \item \textbf{Аномальность (cross-fitted)}: расстояние текущего дня до
  медианного референса из \emph{предыдущих} $\le 3$ дней (без включения
  текущего дня).
\end{itemize}

\section{Модели}
Базовая подтверждающая модель:
\begin{itemize}
  \item иерархическая байесовская логистическая регрессия;
  \item случайные intercepts для крысы и месяца;
  \item регуляризующие нормальные/half-normal priors.
\end{itemize}

Сравниваются четыре варианта:
\begin{enumerate}
  \item \texttt{null}: только глобальный intercept;
  \item \texttt{calendar}: intercept + rat/month random intercepts;
  \item \texttt{full}: календарь + все признаки;
  \item \texttt{no\_shape}: календарь + признаки без shape/anomaly/pairwise.
\end{enumerate}

\section{Валидация}
\begin{itemize}
  \item outer grouped CV по \texttt{group\_event\_date} (все крысы одной даты
  в одном фолде);
  \item отдельный temporal holdout (ранние даты train, поздние test);
  \item метрики: log score, Brier, ROC-AUC, PR-AUC, ECE (10 bins).
\end{itemize}

\section{Где смотреть результаты}
После запуска скриптов:
\begin{itemize}
  \item \texttt{run\_output\_risk\_model/risk\_daily\_dataset.csv},
  \texttt{risk\_daily\_dataset\_summary.json};
  \item \texttt{run\_output\_risk\_model/risk\_model\_cv\_metrics.csv},
  \texttt{risk\_model\_run\_summary.json}.
\end{itemize}

\section{Что ещё нужно для финализации}
\begin{enumerate}
  \item Подключение к БД с интенсивностями для внешних ковариат и стратификаций
  по порогу события.
  \item Предзаданный список quiet-периодов и правила матчирования.
  \item Репликация на независимом временном диапазоне.
\end{enumerate}

\end{document}
